
\subsection{不可定向族指标之例}
\label{subsec:8-1-2}

本小节将说明：若不假设 $F$ 平凡，命题~\ref{prop:8-1-4}
证明中出现的指标丛确实可能不可定向。

可以如下找到这样的丛对 $(E,F)$。取平凡丛
\[
 \widetilde E=(D^2\times[0,1])\times\mathbb C^n
\]
以及边界上的实子丛%
\translatornote{源书把 $\widetilde F$ 误排为
$\mathbb R^n\times(\partial D^2\times[0,1])\times\mathbb C^n$；
这多出的 $\mathbb C^n$ 因子使其不可能是 $\widetilde E$ 的实子丛。
这里依紧接的“实部”表述及所需实秩 $n$ 恢复。}
\[
 \widetilde F=(\partial D^2\times[0,1])\times\mathbb R^n.
\]
识别 $D^2\times\{0\}$ 与 $D^2\times\{1\}$，便得到
$D^2\times S^1$。把这一识别提升到向量丛 $\widetilde E$，使其在
$\partial D^2\times\{i\}$（$i=0,1$）上的限制保持实部
$\widetilde F$，并以一个同伦非平凡回路
\[
 \gamma\colon\partial D^2\longrightarrow SO(n)
\]
识别两端。这样的识别存在，因为该回路在包含映射
$SO(n)\to U(n)$ 下的像是零同伦的，因而可延拓为连续映射
$D^2\to U(n)$。记由此识别所得的丛对为 $(E,F)$。

下面的命题~\ref{prop:8-1-7} 蕴含：对这样的 $(E,F)$，指标丛
不可定向。这意味着，用 $\gamma$ 改变稳定平凡化会逆转全纯圆盘
模空间的定向。

来考察 $n=2$ 的情形。所得 $D^2\times S^1$ 上丛对 $(E,F)$
还有如下描述。取平凡向量丛
\[
 (D^2\times[0,1])\times\mathbb H,
\]
其中四元数 $\mathbb H$ 被视为左 $\mathbb C$-模。以右乘
$i\in\mathbb H^*$ 识别 $D^2\times\{0\}$ 与
$D^2\times\{1\}$ 上的纤维，即规定
\[
 (z,0,\zeta i)\sim(z,1,\zeta)
 \qquad(z\in D^2,\ \zeta\in\mathbb H),
\]
由此得到 $E$。对 $(z,t)\in\partial D^2\times[0,1]$，子空间
\[
 F_{z,t}=(\mathbb R\,1\oplus\mathbb R\,j)
 \left\{
 t\left(\frac{z+\bar z}{2}+\frac{z-\bar z}{2i}j\right)
 +(1-t)k
 \right\}
\]
是 $\mathbb H$ 中的全实子空间。因此，它们组成平凡丛
$(D^2\times[0,1])\times\mathbb H$ 在
$\partial D^2\times[0,1]$ 上的全实子丛，并下降为 $E$ 在
$\partial D^2\times S^1$ 上的全实子丛 $F$。

\begin{proposition}{8.1.7}\label{prop:8-1-7}
由 $(E,F)$ 扭曲的多尔博算子之族指标，是 $S^1$ 上不可定向的
虚向量丛。
\end{proposition}

\begin{proof}
利用傅里叶展开，我们显式计算多尔博算子的核与余核。
取左 $\mathbb C$-模 $\mathbb H$ 的基
$\langle1+k,1-k\rangle$。把 $\mathbb H$ 看成
$\mathbb C\oplus\mathbb Cj$，便可将此基写成列向量
\[
 \left\langle
 \begin{pmatrix}1\\ i\end{pmatrix},
 \begin{pmatrix}1\\ -i\end{pmatrix}
 \right\rangle.
\]
设 $f\colon D^2\to\mathbb C^2=\mathbb H$ 为全纯映射。于是
$f$ 有傅里叶展开
\[
 f(z)=\sum_{n=0}^{\infty}z^n
 \left\{
 \alpha_n\begin{pmatrix}1\\ i\end{pmatrix}
 +\beta_n\begin{pmatrix}1\\ -i\end{pmatrix}
 \right\},
 \qquad z\in D^2,\quad \alpha_n,\beta_n\in\mathbb C.
\]
现假设 $f|_{\partial D^2}$ 的像位于 $\mathbb H$ 的全实子空间
\[
 (\mathbb R\,1\oplus\mathbb R\,j)
 \left\{
 t\left(\frac{z+\bar z}{2}+\frac{z-\bar z}{2i}j\right)
 +(1-t)k
 \right\}.
\]
这个条件等价于
\begin{equation}
 f(z)\left\{
 t\left(\frac{z+\bar z}{2}-\frac{z-\bar z}{2i}j\right)
 -(1-t)k
 \right\}\in\mathbb R\,1\oplus\mathbb R\,j
 \qquad(|z|=1).
 \tag{8.1.8}\label{eq:8-1-8}
\end{equation}

把右乘
\[
 t\left(\frac{z+\bar z}{2}-\frac{z-\bar z}{2i}j\right)
 -(1-t)k
\]
写成从左作用于 $\mathbb H=\mathbb C\oplus\mathbb Cj$ 列向量的
$2\times2$ 复矩阵，便得到
\[
\begin{aligned}
 &t\frac{z+\bar z}{2}
   \begin{pmatrix}1&0\\0&1\end{pmatrix}
 -t\frac{z-\bar z}{2i}
   \begin{pmatrix}0&-1\\1&0\end{pmatrix}
 -(1-t)\begin{pmatrix}0&i\\i&0\end{pmatrix}
 \\
 &\qquad={}
 z\frac t2\begin{pmatrix}1&-i\\i&1\end{pmatrix}
 +\bar z\frac t2\begin{pmatrix}1&i\\-i&1\end{pmatrix}
 -(1-t)\begin{pmatrix}0&i\\i&0\end{pmatrix}.
\end{aligned}
\]
因此
\[
\begin{aligned}
 &f(z)\left\{
 t\left(\frac{z+\bar z}{2}-\frac{z-\bar z}{2i}j\right)
 -(1-t)k
 \right\}
 \\
 ={}&\sum_{n=0}^{\infty}z^n\Biggl\{
 z\frac t2\alpha_n
  \begin{pmatrix}1&-i\\i&1\end{pmatrix}
  \begin{pmatrix}1\\i\end{pmatrix}
 +\bar z\frac t2\alpha_n
  \begin{pmatrix}1&i\\-i&1\end{pmatrix}
  \begin{pmatrix}1\\i\end{pmatrix}
 \\
 &\qquad
 -(1-t)\alpha_n
  \begin{pmatrix}0&i\\i&0\end{pmatrix}
  \begin{pmatrix}1\\i\end{pmatrix}
 +z\frac t2\beta_n
  \begin{pmatrix}1&-i\\i&1\end{pmatrix}
  \begin{pmatrix}1\\-i\end{pmatrix}
 \\
 &\qquad
 +\bar z\frac t2\beta_n
  \begin{pmatrix}1&i\\-i&1\end{pmatrix}
  \begin{pmatrix}1\\-i\end{pmatrix}
 -(1-t)\beta_n
  \begin{pmatrix}0&i\\i&0\end{pmatrix}
  \begin{pmatrix}1\\-i\end{pmatrix}
 \Biggr\}.
\end{aligned}
\]
注意当 $|z|=1$ 时 $\bar z=z^{-1}$，故上式可写成
\[
 \sum_{n=-1}^{\infty}z^n\vec a_n,
 \qquad \vec a_n\in\mathbb C\oplus\mathbb Cj,
\]
其中
\[
\begin{aligned}
 \vec a_{-1}
 &={}
 t\beta_0\begin{pmatrix}1\\-i\end{pmatrix},
 \\
 \vec a_0
 &={}
 -(1-t)\beta_0\begin{pmatrix}1\\i\end{pmatrix}
 +\{(1-t)\alpha_0+t\beta_1\}
  \begin{pmatrix}1\\-i\end{pmatrix},
 \\
 \vec a_n
 &={}
 \{t\alpha_{n-1}-(1-t)\beta_n\}
  \begin{pmatrix}1\\i\end{pmatrix}
 +\{(1-t)\alpha_n+t\beta_{n+1}\}
  \begin{pmatrix}1\\-i\end{pmatrix}
 \qquad(n\geq1).
\end{aligned}
\]
因为 $\mathbb R\,1\oplus\mathbb R\,j$ 是 $\mathbb H$ 中的
全实子空间，条件~\eqref{eq:8-1-8} 蕴含
$\overline{\vec a_n}=\vec a_{-n}$。这些条件给出
\begin{align}
 \overline{\vec a_{-1}}&=\vec a_1,
  \tag{8.1.9.1}\label{eq:8-1-9-1}\\
 \vec a_0&\in\mathbb R^2,
  \tag{8.1.9.2}\label{eq:8-1-9-2}\\
 \vec a_n&=0\qquad(n\geq2).
  \tag{8.1.9.3}\label{eq:8-1-9-3}
\end{align}

先设 $t\neq0,1$。由~\eqref{eq:8-1-9-3}，
\[
 \alpha_n=0\quad(n\geq1),
 \qquad
 \beta_n=0\quad(n\geq2).
\]
条件~\eqref{eq:8-1-9-2} 等价于
\[
 -(1-t)\overline{\beta_0}=(1-t)\alpha_0+t\beta_1.
\]
结合~\eqref{eq:8-1-9-1}，容易看出 $\alpha_0$ 与 $\beta_1$
由 $\beta_0$ 决定；显式地，
\[
 \alpha_0=\frac{2t-1}{2t^2-2t+1}\,\overline{\beta_0},
 \qquad
 \beta_1=-\frac{1-t}{t}
 \bigl(\overline{\beta_0}+\alpha_0\bigr).
\]
所以 $f$ 只由 $\beta_0\in\mathbb C$ 决定。因此，对
$t\neq0,1$，多尔博算子的核之实维数为 $2$。另一方面，该算子的
弗雷德霍姆指标也是 $2$，故该算子满射。

再描述 $t=0,1$ 时的核。在两种情形下，由
\eqref{eq:8-1-9-3} 均有
$\alpha_n=0$（$n\geq1$）和 $\beta_n=0$（$n\geq2$）。此外，
由~\eqref{eq:8-1-9-2}，两种情形下均有 $\beta_1=0$，且
\[
 \alpha_0=
 \begin{cases}
  -\overline{\beta_0},&t=0,\\
  \phantom{-}\overline{\beta_0},&t=1.
 \end{cases}
\]
故 $\beta_0$ 在两种情形下都决定算子的核。

有了这些核的显式描述，就可以证明指标丛不可定向。
对 $t\in[0,1]$，核同构于 $\{\beta_0\in\mathbb C\}$，
因而带有复定向。我们要证明的是，$t=0$ 与 $t=1$ 处两核之间的
识别映射反定向。当 $t=0$ 时，核 $V_0$ 为
\[
 V_0=
 \left\{
 -\overline{\beta_0}\begin{pmatrix}1\\i\end{pmatrix}
 +\beta_0\begin{pmatrix}1\\-i\end{pmatrix}
 \mathrel{\Big|}\beta_0\in\mathbb C
 \right\};
\]
当 $t=1$ 时，核 $V_1$ 为
\[
 V_1=
 \left\{
 \overline{\beta_0}\begin{pmatrix}1\\i\end{pmatrix}
 +\beta_0\begin{pmatrix}1\\-i\end{pmatrix}
 \mathrel{\Big|}\beta_0\in\mathbb C
 \right\}.
\]
回忆端点纤维由右乘 $i$ 的关系识别。仍如上把右乘写成从左作用于
$\mathbb C\oplus\mathbb Cj$ 的 $2\times2$ 复矩阵，即
\[
 \begin{pmatrix}i&0\\0&-i\end{pmatrix}.
\]
于是，在这一识别下，与 $V_0$ 比较的是
\[
\begin{aligned}
 V_1\cdot i
 &={}
 \left\{
 \overline{\beta_0}
  \begin{pmatrix}i&0\\0&-i\end{pmatrix}
  \begin{pmatrix}1\\i\end{pmatrix}
 +\beta_0
  \begin{pmatrix}i&0\\0&-i\end{pmatrix}
  \begin{pmatrix}1\\-i\end{pmatrix}
 \right\}
 \\
 &={}
 \left\{
 \overline{\beta_0}\begin{pmatrix}i\\1\end{pmatrix}
 +\beta_0\begin{pmatrix}i\\-1\end{pmatrix}
 \mathrel{\Big|}\beta_0\in\mathbb C
 \right\}.
\end{aligned}
\]
这个识别映射反定向。事实上，依次取 $\beta_0=1,i$，便分别得到
实向量空间 $V_0$ 与 $V_1\cdot i$ 的基
\[
 \left\langle
  \begin{pmatrix}0\\-2i\end{pmatrix},
  \begin{pmatrix}2i\\0\end{pmatrix}
 \right\rangle,
 \qquad
 \left\langle
  \begin{pmatrix}2i\\0\end{pmatrix},
  \begin{pmatrix}0\\-2i\end{pmatrix}
 \right\rangle.
\]
显然，这两个空间的定向相反。命题~\ref{prop:8-1-7} 得证。
\end{proof}

注意，上述全实子丛 $F$ 可定向，但不是 spin 的。
